\textbf{1. Input data.} The specified dimensions and sound speed are
\[
\boxed{L_{0}=6\,\mathrm{mm}=0.006\,\mathrm{m},\qquad
L_{p}=6\,\mathrm{mm},\qquad L_{w}=6\,\mathrm{mm},\qquad
c=340\,\mathrm{m/s}.}
\]

The ruler markings in Figure~\ref{fig:hw02-04} are used as the reference scale. I measured the overall whistle from the left mouthpiece tip to the right piston-handle tip. This gives the estimates
\[
\boxed{L_{e}\approx 0.210\,\mathrm{m},\qquad L_{c}\approx 0.135\,\mathrm{m}.}
\]
The photograph provides an overall length estimate rather than a direct measurement of the internal air-column endpoints. The fixed piston and wall thicknesses therefore remain input dimensions, while their contributions cancel when the change in overall position is calculated.

\textbf{2. Fully extended whistle length.} Reading the upper ruler in Figure~\ref{fig:hw02-04}(b) from the left mouthpiece tip to the right handle tip gives
\[
\boxed{L_{e}\approx 0.210\,\mathrm{m}.}
\]

\textbf{3. Fully contracted whistle length.} Reading the upper ruler in Figure~\ref{fig:hw02-04}(a) between the same two physical reference points gives
\[
\boxed{L_{c}\approx 0.135\,\mathrm{m}.}
\]

\textbf{4. Fully extended travel.} The fixed wall and piston thicknesses are present in both overall-length measurements. Thus, the slider travel is
\[
\boxed{dL_{\mathrm{I}}=L_{e}-L_{c}=0.210-0.135=0.0750\,\mathrm{m}.}
\]

\textbf{5. Fully extended cavity length.} The problem directs us to recall the closed-closed model from Problem B.3. Using the residual cavity length as the fixed offset gives
\[
\boxed{L_{\mathrm{I}}=L_{0}+dL_{\mathrm{I}}=0.006+0.075=0.0810\,\mathrm{m}.}
\]

\textbf{6. Fully extended frequency.} For the assigned closed-closed model, the fundamental frequency is $f=c/(2L)$. Therefore,
\[
f_{\mathrm{I}}=\frac{340}{2(0.081)}=2098.77\,\mathrm{Hz}\approx 2.10\,\mathrm{kHz}.
\]
The nearest equal-tempered note is
\[
\boxed{f_{\mathrm{I}}\approx2.10\,\mathrm{kHz},\qquad \text{nearest note }\mathrm{C7}\ (2093.00\,\mathrm{Hz}).}
\]

\textbf{7. Forty-percent extension travel.} The travel at $40\%$ extension is
\[
\boxed{dL_{\mathrm{II}}=0.40dL_{\mathrm{I}}=0.40(0.075)=0.0300\,\mathrm{m}.}
\]

\textbf{8. Forty-percent extension cavity length.} Adding the residual cavity gives
\[
\boxed{L_{\mathrm{II}}=L_{0}+dL_{\mathrm{II}}=0.006+0.030=0.0360\,\mathrm{m}.}
\]

\textbf{9. Forty-percent extension frequency.} The closed-closed fundamental is
\[
f_{\mathrm{II}}=\frac{340}{2(0.036)}=4722.22\,\mathrm{Hz}.
\]
The nearest equal-tempered note is
\[
\boxed{f_{\mathrm{II}}\approx4.72\,\mathrm{kHz},\qquad \text{nearest note }\mathrm{D8}\ (4698.64\,\mathrm{Hz}).}
\]

\textbf{10. Comment.} The ruler estimates give a travel of approximately $75\,\mathrm{mm}$, and the model predicts that increasing the cavity length lowers the fundamental frequency. The two frequencies follow this inverse relationship. A physical slide whistle has an open fipple and a closed piston, so the course notes' closed-open model (pp. 114--117) gives $f=c/(4L)$ for that configuration. The corresponding estimates would be $1049.38\,\mathrm{Hz}$ (approximately C6) at full extension and $2361.11\,\mathrm{Hz}$ (approximately D7) at $40\%$ extension. The closed-closed results above follow the assignment instruction to recall Problem B.3; the open-fipple behavior is the relevant physical limitation of that idealization.

\textbf{11. Frequency of B7.} Relative to A4, B7 is 38 semitones higher. Equal temperament gives
\[
\boxed{f_{\mathrm{III}}=440\,2^{38/12}=3951.07\,\mathrm{Hz}\approx3.95\,\mathrm{kHz}.}
\]

\textbf{12. Cavity length for B7.} Applying the assigned closed-closed fundamental relation gives
\[
\boxed{L_{\mathrm{III}}=\frac{c}{2f_{\mathrm{III}}}
=\frac{340}{2(3951.07)}\approx0.0430\,\mathrm{m}.}
\]

\textbf{13. Extension travel for B7.} The required travel beyond the residual cavity is
\[
\boxed{dL_{\mathrm{III}}=L_{\mathrm{III}}-L_{0}=0.04303-0.006\approx0.0370\,\mathrm{m}.}
\]

\textbf{14. Required piston extension.} Relative to the full measured travel, the required percentage is
\[
\boxed{100\frac{dL_{\mathrm{III}}}{dL_{\mathrm{I}}}
=100\frac{0.03703}{0.075}=49.4\%.}
\]

\textbf{15. Comment.} The B7 calculation places the slider at approximately half of the available measured travel under the assigned closed-closed model. The result is within the photographed travel range. Using the physical closed-open fipple model would require $L_{\mathrm{III}}=0.02151\,\mathrm{m}$, or approximately $20.7\%$ of the measured travel. Thus, the boundary-condition model determines the numerical prediction.

\clearpage
\matlabcaption{MATLAB calculations for B.5.}
\lstinputlisting[style=hwmatlab]{matlab/B5.m}
